\documentclass[
  aps,
  prc,
  reprint,
  superscriptaddress,
  floatfix,
  nofootinbib
]{revtex4-2}

% ---------------------------------------------------------------------------
%  DRAFT SKELETON (no calculations yet) -- 2026-10-10
%  Resolution refinement + rodeo algorithm applied to the two-fragment
%  pairing model of the companion draft paper2/twontransfer.pdf.
%  Everything marked \todo{...} is a placeholder to be filled from the
%  notebook once the calculations exist.  Numbers quoted without a \todo
%  are taken from the companion draft or from the first paper
%  (paper/eigenq_pairing.tex) or are pure combinatorics.
% ---------------------------------------------------------------------------

\usepackage{amsmath,amssymb}
\usepackage{physics}
\usepackage{graphicx}
\usepackage{xcolor}
\usepackage{hyperref}

\newcommand{\todo}[1]{\textcolor{red}{[#1]}}
\newcommand{\Hlow}{H_{\mathrm{low}}}
\newcommand{\Hhigh}{H_{\mathrm{high}}}
\newcommand{\Eobj}{E_{\mathrm{obj}}}
\newcommand{\HP}{\hat H_{P}}
\newcommand{\HT}{\hat H_{T}}
\newcommand{\HPT}{\hat H_{PT}}
\newcommand{\Snum}{S_{\mathrm{num}}}
\newcommand{\Scfg}{S_{\mathrm{cfg}}}
\newcommand{\Sdiag}{S_{\mathrm{diag}}}
\newcommand{\Sn}{^{116}\mathrm{Sn}}
\newcommand{\Ni}{^{60}\mathrm{Ni}}
\newcommand{\Zr}{^{96}\mathrm{Zr}}
\newcommand{\Ca}{^{40}\mathrm{Ca}}

\begin{document}

\title{Optimizer-free quantum simulation of sub-barrier neutron-pair transfer:
resolution refinement and rodeo spectroscopy in the two-fragment pairing model}

% \todo{Author list: union of the companion-draft authors and the
% refinement/rodeo paper authors, order to be agreed.}
\author{Morten Hjorth-Jensen}
\email{mhjensen@uio.no}
\affiliation{Department of Physics, University of Oslo, N-0316 Oslo, Norway}
\author{\todo{Co-authors}}
\affiliation{\todo{Affiliations}}
\date{\today}

\begin{abstract}
The number-conserving two-fragment pairing model of
Ref.~\cite{TwoFragment2026} reduces sub-barrier neutron-pair transfer between
superfluid nuclei to a time-dependent quasi-spin Hamiltonian that maps
exactly onto longitudinal fields and number-conserving $XY$ couplings on one
qubit per pair slot.  Running that model on a quantum device requires three
things that the classical benchmark of Ref.~\cite{TwoFragment2026} supplies
for free: the correlated ground states of the two isolated fragments as the
entrance state, a digital propagation through the collision pulse, and a
readout that resolves the exit channel in pair number \emph{and} in
fragment excitation energy.  We show how the first and the third of these
can be done without any variational optimization by the two primitives that
we implemented and assessed on the constant-pairing Hamiltonian in
Ref.~\cite{EigenQPairing2026}: \emph{resolution refinement}, the adiabatic
transport of an eigenstate from a small to a larger active orbital space,
and the \emph{rodeo algorithm}, a projective spectral filter.  Each fragment
ground state is grown orbital by orbital from a single filled pair-slot
configuration, so the entrance state of $\Sn+\Ni$ (and the exit-channel
reference states) are prepared with zero classical parameters; the same
Trotterized $XY$ evolution that implements the refinement path implements the
collision; and a rodeo filter run with the two commuting fragment
Hamiltonians after the pulse measures the exclusive, state-resolved transfer
probabilities $P_{\kappa ab}$ from which the inclusive $P_{2n}$ compared with
experiment, the ground-state-to-ground-state channel, and a rigorous
two-sided bracket $\Snum\le S_{T|P}\le\Sdiag$ on the interfragment
entanglement entropy follow.  \todo{Results sentence(s): refinement times,
fidelities, gate budgets for 11, 19 and 30 qubits; comparison of the
projector and the orbital-lowering refinement paths; rodeo resolution of
the $0^+$ strength in $^{42}$Ca for the inclusive $\Zr+\Ca$ data.}  All
stages are verified against the exact finite-space propagation of
Ref.~\cite{TwoFragment2026}.
\end{abstract}

\maketitle

% ===========================================================================
\section{Introduction}
% ===========================================================================
\todo{Paragraph 1: pairing, pair transfer, sub-barrier regime; one paragraph
condensed from the companion draft.}

\todo{Paragraph 2: why quantum simulation of transfer; the two-fragment
pairing model as the compact, number-conserving, classically verifiable
benchmark; 30 qubits for $\Sn+\Ni$, 26 for $\Zr+\Ca$; the model reproduces
the measured $D$ dependence with two calibrated parameters.}

Paragraph 3 (what is missing between the model and a device).  The
classical benchmark of Ref.~\cite{TwoFragment2026} starts from the exact
ground states of the isolated fragments, propagates every amplitude in the
5554-dimensional seniority-zero space, and reads off the complete transfer
distribution and the reduced density matrix of the target.  A device has
none of this.  It must (i) \emph{prepare} the correlated entrance state
$\ket{\mathrm{GS}_T(2)}\otimes\ket{\mathrm{GS}_P(5)}$ on 11 and 19 qubits;
(ii) \emph{propagate} it with a product formula through the Gaussian
coupling pulse; and (iii) \emph{read out} more than the pair-number
partition, because the quantities that carry the physics --- the
ground-state-to-ground-state channel compared with exclusive data, the
excited $0^+$ strength lumped into inclusive data, and the entanglement
entropy --- live in the energy distribution within each transfer sector.

Paragraph 4 (what we bring from Ref.~\cite{EigenQPairing2026}).  In
Ref.~\cite{EigenQPairing2026} we implemented, on the constant-pairing
Hamiltonian, a pipeline in which an eigenstate is transported adiabatically
from a coarse to a fine single-particle basis along a shift-controlled,
gapped path (resolution refinement~\cite{Bogner2026}) and then cleaned by
the rodeo algorithm~\cite{Choi2021,Qian2024}.  The assessment was sober:
at eight qubits refinement did not beat a cheap UCCSD input on gate count;
its merits were the absence of a classical optimization loop and of
training measurements, and the fact that it is built from the same
Trotterized evolution a device runs anyway.  The two-fragment transfer
model is a setting in which exactly those merits matter: the ``coarse''
states are \emph{single computational basis states} (filled pair slots), the
refinement, the collision and the rodeo controlled evolutions are all
exponentials of one and the same two-local $XY$ Hamiltonian, and the
observable of interest is not a ground-state energy but a distribution over
product eigenstates, which is precisely what a spectral filter measures.

Paragraph 5 (contributions and outline).  \todo{List: (a) orbital-by-orbital
refinement chains for $\Ni$ and $\Sn$ from a filled configuration, two path
variants (projector path of Ref.~\cite{EigenQPairing2026} and a two-local
``orbital-lowering'' path); (b) rodeo filtering of the refined fragment
states; (c) the collision as the same Trotterized primitive, with an
interaction-picture split that removes the large single-particle energies
from the Trotter error; (d) two-fragment rodeo spectroscopy of the exit
channel and the entropy bracket; (e) gate budgets per accepted preparation,
collision and readout, verified against exact evolution.  Outline of
sections.}

% ===========================================================================
\section{The two-fragment pairing model}
\label{sec:model}
% ===========================================================================
We summarize the model of Ref.~\cite{TwoFragment2026} in the notation used
there.  Each fragment $X\in\{P,T\}$ is a set of active neutron orbitals $p$
with single-particle energy $\epsilon_{pX}$ and pair degeneracy
$\Omega_{pX}=(2j_{pX}+1)/2$.  Restricted to seniority zero, the isolated
fragment is the constant-strength pairing Hamiltonian
\begin{equation}
  \hat H_X=\sum_{p\in X}2\epsilon_{pX}\hat n_{pX}
  -G_X\sum_{p,p'\in X}\hat P^\dagger_{pX}\hat P_{p'X},
  \qquad G_X=\frac{20.5}{A_X}\,\mathrm{MeV},
  \label{eq:HX}
\end{equation}
with $M_X$ active pairs, and the collision is
\begin{equation}
  \hat H(t,D)=\HP+\HT+\HPT(t,D),\quad
  \HPT=-G_{PT}(t,D)\sum_{p\in P,\,q\in T}
  \big(\hat P^\dagger_{qT}\hat P_{pP}+\mathrm{h.c.}\big),
  \label{eq:Hcoll}
\end{equation}
with a Gaussian pulse $G_{PT}(t,D)=G_{\rm pk}(D)e^{-(t/\tau(D))^2}$ and
$\tau(D)$ fixed by the Coulomb curvature.  For $\Sn+\Ni$ the active spaces
are those of Table~I of Ref.~\cite{TwoFragment2026}: $\Ni$ has
$\Omega=(2,1,3,5)$ for $(2p_{3/2},2p_{1/2},1f_{5/2},1g_{9/2})$ with $M_T=2$,
and $\Sn$ has $\Omega=(4,1,2,6,4,2)$ for
$(1g_{7/2},3s_{1/2},2d_{3/2},1h_{11/2},2f_{7/2},3p_{3/2})$ with $M_P=5$; the
fixed-number dimensions are $d_T(2)=9$, $d_P(5)=154$, and the $M=7$ collision
space has dimension 5554.

\paragraph{Pair-slot encoding.}
Resolving every orbital into $\Omega_{pX}$ hard-core pair slots
$i=(X,p,\lambda)$, with $\hat P^\dagger_{pX}=\sum_\lambda\hat b^\dagger_{i}$,
is an exact representation of the seniority-zero space on the
permutation-symmetric subspace of each orbital
(Appendix~B of Ref.~\cite{TwoFragment2026}), and gives
\begin{equation}
  \hat H(t,D)=E_c+\sum_i\zeta_iZ_i+\sum_{i<j}J_{ij}(t,D)\,(X_iX_j+Y_iY_j),
  \label{eq:HXY}
\end{equation}
with $\zeta_i=-\epsilon_i+G_{ii}/2$ and $J_{ij}=-G_{ij}/2$, where
$G_{ij}=G_X$ inside a fragment and $G_{PT}(t,D)$ across.  For $\Sn+\Ni$ this
is $N_q=11+19=30$ qubits, seven of which are occupied, and the Hamiltonian
has 30 single-qubit $Z$ terms and $\binom{11}{2}+\binom{19}{2}+11\cdot19
=55+171+209=435$ two-local $XY$ terms, of which the 209 interfragment
terms are switched on only during the pulse.  Two features of
Eq.~\eqref{eq:HXY} drive everything below.  First, the encoding is
\emph{two-local}: no Jordan--Wigner strings appear, in contrast to the
eight-qubit pairing model of Ref.~\cite{EigenQPairing2026} where a Trotter
step of $k=4$ levels required 325 Pauli strings of weight up to eight.
Second, the total pair number $\hat M=\sum_i\hat n_i$ and each fragment's
pair number $\hat M_X$ (outside the pulse) are conserved, so every stage
stays in a fixed-$M$ sector and the sector partition $\kappa$ is read out by
$Z$ measurements on the target register alone.

\paragraph{Observables.}
After the pulse the state is analyzed in the product eigenbasis of
$\HP+\HT$,
\begin{equation}
  \ket{\Psi(t_f,D)}=\sum_{\kappa}\sum_{a,b}c_{\kappa ab}(D)\,
  \ket{E^{a}_{T}(M_T{+}\kappa)}\otimes\ket{E^{b}_{P}(M_P{-}\kappa)},
  \label{eq:final}
\end{equation}
where $a,b=0$ label the fragment ground states.  Ref.~\cite{TwoFragment2026}
compares $P_{+1}/P_0$ with $P_\kappa=\sum_{ab}\abs{c_{\kappa ab}}^2$ to the
measured $P_{2n}(D)$, and splits the target entanglement entropy as
$S_{T|P}=\Snum+\Scfg$ with $\Snum=-\sum_\kappa P_\kappa\log_2P_\kappa$.  At
the reference distance $D=13.4669$~fm the exact values are $P_0=0.97209$,
$P_{+1}=0.02206$, $P_{-1}=0.00566$, $P_{+2}=1.96\times10^{-4}$ and
$S_{T|P}=0.20622$ bits of which $\Scfg=4.31\times10^{-4}$ bits.  These are
the targets every stage below is checked against.

% ===========================================================================
\section{The optimizer-free pipeline}
\label{sec:pipeline}
% ===========================================================================
Figure~\ref{fig:pipeline} summarizes the three stages: entrance-state
preparation by refinement and rodeo (Secs.~\ref{sec:refine},
\ref{sec:rodeo-prep}), the Trotterized collision (Sec.~\ref{sec:collision}),
and exit-channel rodeo spectroscopy (Sec.~\ref{sec:rodeo-exit}).  All three
are exponentials of Eq.~\eqref{eq:HXY} with different coefficient schedules.

\begin{figure}[t]
  \centering
  \fbox{\parbox{0.95\columnwidth}{\centering\vspace{2.5cm}\todo{Fig.~1:
  pipeline schematic.  Left: two registers (11 target slots, 19 projectile
  slots) grown orbital by orbital from filled configurations; rodeo ancilla.
  Middle: Gaussian pulse $G_{PT}(t)$, interfragment $XY$ terms switched on.
  Right: $Z$ readout of $\hat M_T$ plus two-fragment rodeo scan in
  $(E_T,E_P)$.}\vspace{2.5cm}}}
  \caption{Schematic of the optimizer-free pipeline.  \todo{caption}}
  \label{fig:pipeline}
\end{figure}

% ---------------------------------------------------------------------------
\subsection{Resolution refinement of the fragment ground states}
\label{sec:refine}
% ---------------------------------------------------------------------------
In Ref.~\cite{EigenQPairing2026} ``resolution'' meant the number $k$ of
doubly degenerate levels at fixed particle number.  Here it means the set of
active orbitals of a fragment at fixed pair number $M_X$, which is exactly
the ``larger active spaces'' extension anticipated in the conclusions of
Ref.~\cite{TwoFragment2026}, and, read backwards, the frozen-orbital
approximation being released adiabatically.  Choose a chain of orbital sets
\begin{equation}
  \mathcal A_0\subset\mathcal A_1\subset\cdots\subset\mathcal A_R=X,
\end{equation}
and let $\mathcal S_r$ be the slots of $\mathcal A_r$.  The prolongation
$P_r$ embeds a state on $\mathcal S_r$ into $\mathcal S_{r+1}$ by appending
the new slots in $\ket{0}$; in the pair-slot encoding this is gate-free, as
it was in canonical Jordan--Wigner order in Ref.~\cite{EigenQPairing2026},
and the prolonged state satisfies
$\expval{\hat H_{\mathcal A_{r+1}}}{P_r\psi_r}=\expval{\hat H_{\mathcal A_r}}{\psi_r}$
exactly.

\paragraph{Chains.}
The occupations of Table~I of Ref.~\cite{TwoFragment2026} suggest the chains
of Table~\ref{tab:chains}.  Both start from a \emph{filled} configuration,
i.e.\ a single computational basis state: $\Ni$ with two pairs in the two
slots of $2p_{3/2}$, and $\Sn$ with five pairs in the $4+1$ slots of
$1g_{7/2}$ and $3s_{1/2}$.  No variational or classical preprocessing of any
kind is needed to start the chain, which is the analogue of the
$k=2\to4$ bootstrap of Ref.~\cite{EigenQPairing2026} where the $k=2$ coarse
state was the single filled determinant.  \todo{The exact overlaps of the
filled states with the full ground states are to be computed; the mean
occupations $\langle\hat n\rangle/\Omega$ of $3s_{1/2}$ (0.48) and
$2d_{3/2}$ (0.29) in $\Sn$ suggest that the first projectile step carries
most of the correlation, whereas for $\Ni$ the depletion of $2p_{3/2}$
(1.79 of 2 pairs) is mild.}

\begin{table}[t]
  \caption{Candidate refinement chains for the $\Sn+\Ni$ entrance channel.
  $|\mathcal S_r|$ is the number of pair slots (qubits), $d_r$ the dimension
  of the seniority-zero (permutation-symmetric) space at fixed $M_X$,
  $\binom{|\mathcal S_r|}{M_X}$ the dimension of the fixed-number slot
  sector in which a gate-level simulation lives, and $n_{XY}$ the number of
  $XY$ terms of $\hat H_{\mathcal A_r}$.  Overlaps $p_r$ of the prolonged
  state $P_{r-1}\ket{\mathrm{GS}_{r-1}}$ with $\ket{\mathrm{GS}_r}$, minimum
  path gaps and extracted refinement times $T^*_r$ are \todo{to be
  computed}.}
  \label{tab:chains}
  \begin{ruledtabular}
  \begin{tabular}{llcccccc}
  $X$ & $\mathcal A_r$ & $|\mathcal S_r|$ & $d_r$ & $\binom{|\mathcal S_r|}{M_X}$ & $n_{XY}$ & $p_r$ & $T^*_r$\\
  \hline
  $T$ & $\{2p_{3/2}\}$                      & 2  & 1   & 1     & 1   & --- & --- \\
      & $+\{2p_{1/2},1f_{5/2}\}$            & 6  & 5   & 15    & 15  & \todo{} & \todo{} \\
      & $+\{1g_{9/2}\}$                     & 11 & 9   & 55    & 55  & \todo{} & \todo{} \\
  \hline
  $P$ & $\{1g_{7/2},3s_{1/2}\}$             & 5  & 1   & 1     & 10  & --- & --- \\
      & $+\{2d_{3/2}\}$                     & 7  & 5   & 21    & 21  & \todo{} & \todo{} \\
      & $+\{1h_{11/2}\}$                    & 13 & 26  & 1287  & 78  & \todo{} & \todo{} \\
      & $+\{2f_{7/2},3p_{3/2}\}$            & 19 & 154 & 11628 & 171 & \todo{} & \todo{} \\
  \end{tabular}
  \end{ruledtabular}
\end{table}

\paragraph{Path A: the projector path of Ref.~\cite{EigenQPairing2026}.}
Each step follows
\begin{equation}
  H_r(s)=\cos^2\!\big(\tfrac{\pi s}{2}\big)\,P_r\big(\hat H_{\mathcal A_r}-\mu_r\big)P_r^\dagger
        +\sin^2\!\big(\tfrac{\pi s}{2}\big)\,\big(\hat H_{\mathcal A_{r+1}}-\mu_{r+1}\big),
  \label{eq:pathA}
\end{equation}
with the shifts $\mu$ controlling the minimum gap as in
Ref.~\cite{EigenQPairing2026}.  Because the new slots are empty in the
coarse space, the coarse term is the fine Hamiltonian projected onto the
empty-new-slot subspace,
$P_r\hat H_{\mathcal A_r}P_r^\dagger=\Pi_0\hat H_{\mathcal A_{r+1}}\Pi_0$
with $\Pi_0=\prod_{i\in\mathcal S_{r+1}\setminus\mathcal S_r}\ket{0_i}\!\bra{0_i}$.
This is the term whose circuit cost dominated the $3\to4$ chain in
Ref.~\cite{EigenQPairing2026} (a multi-controlled phase on the new qubits,
$\approx350$ CNOTs per application for two new qubits); here a step adds up
to six new slots ($1h_{11/2}$), so the multi-controlled structure is
correspondingly deeper and Path~A is expected to be the expensive option.
\todo{Give the exact decomposition and CNOT count of $\Pi_0 H\Pi_0$ for the
$XY$ form: every $XY$ term and $Z$ term acquires the control string
$\Pi_0$.}

\paragraph{Path B: orbital lowering (a two-local path).}
The same end points are connected by keeping the full two-local Hamiltonian
and sweeping the single-particle energies of the new orbitals down from a
large value,
\begin{equation}
  H^{\Lambda}_r(s)=\hat H_{\mathcal A_{r+1}}
  +\Lambda\,\cos^2\!\big(\tfrac{\pi s}{2}\big)\sum_{i\in\mathcal S_{r+1}\setminus\mathcal S_r}\hat n_i ,
  \label{eq:pathB}
\end{equation}
i.e.\ $\epsilon_{p}(s)=\epsilon_p+\tfrac12\Lambda\cos^2(\pi s/2)$ for the new
orbitals.  At $s=0$ and $\Lambda\gg G_X\Omega$ the ground state of
Eq.~\eqref{eq:pathB} is the prolonged coarse state up to corrections of
order $(G_X\Omega/\Lambda)^2$ (second-order perturbation theory in the
pair-scattering into the penalized slots eliminates them and reproduces
$\Pi_0\hat H\Pi_0$ to leading order), so Path~B is a two-local relaxation
of Path~A.  Along Path~B every Trotter step costs exactly what a collision
step costs, and the path has a direct physical reading: the frozen orbitals
of the fragment are lowered adiabatically from infinity to their physical
position.  \todo{Compute gap, overlap and $T^*$ along both paths for every
chain step; the comparison of $T^*\times$(cost per step) between A and B is
one of the new results of this paper.  Check the $\Lambda$ dependence of
the initial infidelity and of $T^*$; a schedule with $\Lambda\to\infty$
at $s=0$, e.g.\ $\Lambda\cot^2(\pi s/2)$ regularized, may be preferable.}

\paragraph{Refinement time and sufficient conditions.}
We extract $T^*$ as in Ref.~\cite{EigenQPairing2026}: the smallest $T$ at
which the overlap of the transported state with the exact fine ground state
has reached a stated target, with the $\sin^2$ schedule and the
Jansen--Ruskai--Seiler bound~\cite{Jansen2007} as the reference whose
over-estimation we again quantify.  New here is that the fragment
Hamiltonians are Richardson-type with realistic, non-uniform level spacings
and degeneracies $\Omega_p>1$; the ratios $G_X/\Delta\epsilon$ are
$\approx0.1$--$0.3$ ($G_{\Sn}=0.177$~MeV, $G_{\Ni}=0.342$~MeV), i.e.\ the
fragments are on the weak-coupling side of the $g=1$, $\delta=1$ model of
Ref.~\cite{EigenQPairing2026}, and the minimum path gaps are set by the
lowest seniority-zero excitation of each fragment, which in turn is set
by the orbital spacings rather than by the pairing.  Times are in units of
$\hbar/\mathrm{MeV}=6.58\times10^{-22}$~s; the collision pulse lasts
$10\tau=3.42\,\hbar/$MeV for $\Sn+\Ni$, which is the natural yardstick
for $T^*$.  \todo{Table or figure of spectra along each path; discuss
whether the orbital-spacing-dominated gaps make refinement \emph{faster}
than in the $g=1$ model.}

% ---------------------------------------------------------------------------
\subsection{Rodeo filtering of the refined fragment states}
\label{sec:rodeo-prep}
% ---------------------------------------------------------------------------
Each refined fragment state is cleaned by the rodeo algorithm with
$\Eobj=E^{(0)}_X(M_X)$, exactly as in Sec.~III.D of
Ref.~\cite{EigenQPairing2026}: $M$ cycles of a Hadamard, the controlled
evolution $e^{-i\hat H_Xt_k}$, the phase $e^{i\Eobj t_k}$ and a measurement
of a single reused ancilla; the ensemble infidelity obeys
\begin{equation}
  1-\mathcal F_M=\frac{\sum_{k\neq0}\abs{c_k}^2q_k^M}{p+\sum_{k\neq0}\abs{c_k}^2q_k^M},
  \qquad q_k=\tfrac12\big(1+e^{-(E_k-E_0)^2\sigma^2/2}\big).
  \label{eq:fid}
\end{equation}
Three things differ from the first paper.  (i) The two fragments are
uncoupled at $t_i$, so the entrance state is prepared as two independent
problems on 11 and 19 qubits, each with its own $\Eobj$, and the ancilla can
be shared.  (ii) $\Eobj$ is known classically for the benchmark
($d_P(5)=154$), but a rodeo \emph{scan} on the refined state locates it on a
device as well and doubles as a verification that the refinement reached the
right state --- the energy-scan mode of the algorithm that
Ref.~\cite{EigenQPairing2026} used only for illustration.  (iii) The
controlled evolution is two-local: a controlled $XY$ rotation costs a fixed
small number of CNOTs, and the controlled-reversal-gate construction of
Ref.~\cite{BeeLindgren2022} applies directly to $Z$ and $XY$ terms.  The
width $\sigma$ is chosen from the fragment spectrum so that
$\Delta_{\min}\sigma\gtrsim5$ as before; \todo{quote $\Delta_{\min}$ of
each fragment at fixed $M_X$ once computed}.

\todo{Decide whether to keep a variational comparator.  The natural one in
the $XY$ encoding is pair-UCCD (pUCCD), a product of $XY$ rotations acting
on the filled configuration; it is cheap and would play the role UCCSD
played in Ref.~\cite{EigenQPairing2026}.  Given the main finding there,
the honest framing is: report the matched-observable comparison, but the
selling point of the present pipeline is that no optimizer is needed
anywhere between the filled configuration and the final readout.}

% ---------------------------------------------------------------------------
\subsection{The collision as the same Trotterized primitive}
\label{sec:collision}
% ---------------------------------------------------------------------------
The pulse is propagated with a second-order product formula in the midpoint
rule for the time dependence of $G_{PT}$,
\begin{equation}
  U(t_f,t_i)\approx\prod_{n}
  e^{-i\hat H_Z\delta t/2}\,e^{-i\hat H_{XY}(\bar t_n)\delta t}\,e^{-i\hat H_Z\delta t/2},
  \label{eq:trotter}
\end{equation}
where $\hat H_Z$ collects the 30 longitudinal terms (which commute and are
exponentiated exactly by single-qubit rotations) and $\hat H_{XY}(t)$ the 435
exchange terms, grouped into layers of commuting pairs.  Two observations
specific to this model.

\paragraph{Interaction picture with respect to $\hat H_Z$.}
The single-particle energies are $\sim10$~MeV while the couplings are
$\lesssim0.35$~MeV, so the Trotter error of Eq.~\eqref{eq:trotter} is
dominated by $[\hat H_Z,\hat H_{XY}]$.  Since $\hat H_Z$ is a sum of
commuting single-qubit terms, the rotating frame
$\tilde H_{XY}(t)=e^{i\hat H_Zt}\hat H_{XY}(t)e^{-i\hat H_Zt}$ is still
two-local --- each $X_iX_j+Y_iY_j$ becomes
$\cos(\Delta_{ij}t)(X_iX_j+Y_iY_j)+\sin(\Delta_{ij}t)(X_iY_j-Y_iX_j)$ with
$\Delta_{ij}=2(\zeta_i-\zeta_j)$ --- and its Trotter error is governed by
$G^2$ rather than $G\epsilon$.  Over the pulse duration $3.4\,\hbar/$MeV
and with $G_{\rm pk}\le0.3$~MeV this suggests that a modest number of steps
suffices; \todo{quantify: $P_\kappa(D)$ and $S_{T|P}$ versus $\delta t$ in
both frames against the DOP853 reference of Ref.~\cite{TwoFragment2026};
report the step count at which the three-order-of-magnitude fall of
$P_{2n}(D)$ in Fig.~2 of that paper is reproduced to within the
experimental scatter}.  The same observation applies to the refinement
paths, whose Trotterization in Ref.~\cite{EigenQPairing2026} was found to
be far more forgiving than the sufficient conditions suggest.

\paragraph{Trotter leakage out of the symmetric subspace.}
The exact evolution preserves the permutation-symmetric subspace of every
orbital, but an individual $XY$ term between two slots of the same orbital
does not, so a product formula leaks weight at $O(\delta t^2)$ into
configurations that are not seniority-zero states of the model.  These
configurations have different pairing energies (the diagonal term
$G_Xn(\Omega-n+1)$ is specific to the symmetric irrep), so the leakage is
energetically resolvable: it shows up as extra weight in a rodeo scan and
is removed by the exit-channel filter.  \todo{Measure the leaked weight
versus $\delta t$; it is a cheap and specific Trotter diagnostic for
this encoding.  Consider ordering the within-orbital terms so that the
collective exponential $e^{-i\theta\hat P^\dagger_p\hat P_p}$ is applied as
one block.}

\paragraph{Pulse truncation.}
Ref.~\cite{TwoFragment2026} integrates over $\pm5\tau$ where the coupling
is $e^{-25}$ of its peak.  On a device one truncates to $\pm3\tau$
($e^{-9}\approx10^{-4}$) or less; \todo{report the effect on $P_\kappa$}.

\paragraph{An adiabatic bound on the non-elastic probability.}
The pulse returns $\hat H(t,D)$ to $\HP+\HT$, so the entrance state is the
instantaneous ground state at both ends and $1-P_{000}$, the probability of
having left it, is bounded by the adiabatic theorem with the pulse as the
schedule --- the same Jansen--Ruskai--Seiler expression we evaluated for
the refinement path in Sec.~IV.E of Ref.~\cite{EigenQPairing2026}.  Since
those bounds overestimated the actual leakage by two to three orders of
magnitude, we expect the bound on $1-P_{000}\approx0.03$ to be loose;
\todo{evaluate it anyway, as a link between the two papers and as a
statement about how non-adiabatic sub-barrier transfer is in this model}.

% ---------------------------------------------------------------------------
\subsection{Rodeo spectroscopy of the exit channel}
\label{sec:rodeo-exit}
% ---------------------------------------------------------------------------
This is the stage where the rodeo algorithm does something the first paper
did not ask of it.  After the pulse $\hat H=\HP+\HT$ with $[\HP,\HT]=0$, so
the two fragment Hamiltonians can be filtered \emph{jointly}: $M_T$ rodeo
cycles with $e^{-i\HT t_k}$ and energy $E_T$, followed by $M_P$ cycles with
$e^{-i\HP t_k'}$ and energy $E_P$, with one shared ancilla or two.  For the
state of Eq.~\eqref{eq:final} the all-zero acceptance probability averaged
over Gaussian time sets is
\begin{equation}
  P_{0^{M}}(E_T,E_P)=\sum_{\kappa ab}\abs{c_{\kappa ab}}^2\,
  q^{M_T}\!\big(E^{a}_{T}(M_T{+}\kappa)-E_T\big)\,
  q^{M_P}\!\big(E^{b}_{P}(M_P{-}\kappa)-E_P\big),
  \label{eq:scan2d}
\end{equation}
with $q(\Delta)=\tfrac12(1+e^{-\Delta^2\sigma^2/2})$.  A scan in the
$(E_T,E_P)$ plane therefore displays the final state as peaks at the
product energies, with heights that converge geometrically to the
exclusive probabilities $\abs{c_{\kappa ab}}^2$.  Reading out $\hat M_T$ on
the target register in the same shot assigns each peak its sector $\kappa$
and removes accidental degeneracies between products belonging to
different sectors.  Three uses follow.

\paragraph{Exclusive versus inclusive transfer.}
The quantity compared with the $\Sn+\Ni$ data in Ref.~\cite{TwoFragment2026}
is $P_{+1}/P_0$, the full one-pair-pickup sector.  The filter with
$E_T=E^{(0)}_T(3)$, $E_P=E^{(0)}_P(4)$ measures the ground-state-to-ground-state
channel $\abs{c_{+1,00}}^2$ alone, and the residual
$P_{+1}-\abs{c_{+1,00}}^2$ is the excited $0^+$ strength of
$^{62}$Ni$\otimes{}^{114}$Sn populated in the model.  For $\Zr+\Ca$ the
situation is reversed: the data are inclusive and
Ref.~\cite{TwoFragment2026} had to place the centroid of the $+1$ sector at
$Q=0$ by hand; the two-fragment scan resolves the $0^+$ strength in
$^{42}$Ca orbital by orbital and makes the inclusive-channel approximation
explicit.  Note that the entrance space of $\Ca$ is empty, so for this
reaction only $\Zr$ (three pairs in five orbitals, \todo{slot count})
needs to be prepared, and the target register starts in
$\ket{0\cdots0}$: $\Zr+\Ca$ is the simplest complete end-to-end case and
the natural first calculation.

\paragraph{Measuring a small-weight channel.}
The filter does not need to \emph{prepare} the transferred state with high
fidelity to \emph{measure} its weight: the acceptance probability obeys
$P_M=p+\sum_{k\neq0}\abs{c_k}^2q_k^M$ and $p$ is read off the plateau in
$M$.  For $p=\abs{c_{+1,00}}^2\approx0.02$ the background $2^{-M}(1-p)$
falls below $p$ at $M\approx6$ and below $10^{-2}p$ at $M\approx13$; for the
two-pair channel $P_{+2}\approx2\times10^{-4}$ one needs $M\approx19$.
These cycle counts, times the controlled-evolution cost of
Sec.~\ref{sec:rodeo-prep}, set the readout budget; the acceptance rate is of
course $\approx p$ itself, so the shot overhead is $1/p$, as for any
measurement of a rare channel.  \todo{Table of $M$, $\sigma$, acceptance
and CNOT count per accepted shot for $\kappa=0,\pm1,+2$.}

\paragraph{A two-sided bracket on the entanglement entropy.}
Because $[\rho_T,\hat M_T]=0$, Ref.~\cite{TwoFragment2026} writes
$S_{T|P}=\Snum+\Scfg$ with $\Scfg\ge0$, so $\Snum\le S_{T|P}$; $\Snum$ needs
only the $Z$ readout of $\hat M_T$.  On the other side, the Shannon entropy
of the diagonal of $\rho_T$ in any orthonormal basis is an upper bound on
its von Neumann entropy, and the rodeo scan with $\HT$ alone measures
exactly that diagonal in the eigenbasis of $\HT$,
$P_{\kappa a}=\sum_b\abs{c_{\kappa ab}}^2=\mel{E^a_T}{\rho_T}{E^a_T}$.
Hence
\begin{equation}
  \Snum=-\sum_\kappa P_\kappa\log_2P_\kappa\;\le\;S_{T|P}\;\le\;
  \Sdiag=-\sum_{\kappa,a}P_{\kappa a}\log_2P_{\kappa a},
  \label{eq:bracket}
\end{equation}
and both ends are measured by the same circuit.  Since
Ref.~\cite{TwoFragment2026} finds $\Scfg/S_{T|P}<0.3\%$ at the reference
distance and $\le2\%$ at the smallest, the bracket is expected to be tight
across the whole $D$ range; where it is not, its width is itself a measure
of the coherence within a transfer sector that the diagonal distribution
misses.  \todo{Compute $\Sdiag(D)$ from the exact final states and plot the
bracket against $S_{T|P}(D)$ for both reactions.  The bracket is, to our
knowledge, the first proposal for a device-level estimate of the
interfragment entanglement of Ref.~\cite{TwoFragment2026}; check the
literature on diagonal/``observational'' entropy bounds and cite.}

% ---------------------------------------------------------------------------
\subsection{Optional: the static di-nuclear ground state}
\label{sec:static}
% ---------------------------------------------------------------------------
\todo{Freeze the coupling at its peak value $G_{\rm pk}(D)$ and ask for the
ground state of the coupled 30-qubit Hamiltonian.  The refinement path
$H(s)=\HP+\HT+s\,G_{\rm pk}\hat V_{PT}$ connects the product entrance state to
it (``coupling refinement'', $P=\mathbb 1$, two-local throughout), and the
rodeo filter finishes it.  Its interfragment entropy and pair-number
distribution are the static counterparts of the dynamical ones and give an
adiabatic reference for the sub-barrier transfer.  Decide whether this
belongs in this paper or in a follow-up; it is cheap once the pipeline
exists.}

% ===========================================================================
\section{What carries over from the pairing-model study}
\label{sec:inherit}
% ===========================================================================
\begin{table*}[t]
  \caption{Ingredients of Ref.~\cite{EigenQPairing2026} and their status in
  the present work.}
  \label{tab:inherit}
  \begin{ruledtabular}
  \begin{tabular}{p{0.3\textwidth}p{0.32\textwidth}p{0.32\textwidth}}
  Ingredient & In Ref.~\cite{EigenQPairing2026} & Here \\
  \hline
  Model & constant pairing, $k$ equispaced doubly degenerate levels, $N=4$, $g=1$; optional particle-hole term & two Richardson-type fragments with realistic $\epsilon_p$, $\Omega_p\le6$, $M_P=5$, $M_T=2$; time-dependent interfragment coupling \\
  Encoding & Jordan--Wigner, $2k$ qubits, Pauli strings of weight up to $2k$ & pair slots, one qubit per slot, strictly two-local $Z$/$XY$ \\
  Coarse state & UCCSD--VQE or the filled $k=2$ determinant & filled slot configuration (no parameters) \\
  Prolongation & gate-free in JW order & gate-free \\
  Refinement path & shift-controlled projector path, Eq.~\eqref{eq:pathA} & Path~A unchanged; Path~B (orbital lowering) new \\
  $T^*$ protocol, JRS bound & Sec.~III.C, IV.E & reused unchanged \\
  Rodeo filter & ground-state preparation, Eqs.~\eqref{eq:fid} & reused for fragments; new use as two-fragment spectroscopy, Eq.~\eqref{eq:scan2d} \\
  Comparators & prolonged, UCCSD, refined inputs to the same filter & \todo{pUCCD?}; main comparison is Path~A vs Path~B and prep/collision/readout budgets \\
  Cost accounting & CNOTs per accepted preparation, coarse term included & same, plus collision and readout per accepted shot \\
  Verification & Richardson exact solution, exact sector evolution & exact fragment diagonalization and DOP853 propagation of Ref.~\cite{TwoFragment2026} \\
  Code & \texttt{pairinglib}, notebook, 26 tests & \todo{generalize to $\Omega_p>1$, non-uniform $\epsilon_p$, two registers, time-dependent $J_{ij}$; fixed-$M$ slot-sector simulator ($\binom{30}{7}=2.0\times10^6$ amplitudes for the full collision)} \\
  \end{tabular}
  \end{ruledtabular}
\end{table*}

\todo{Short prose around Table~\ref{tab:inherit}.  Emphasize the two
lessons of the first paper that shaped the design here: (1) the coarse
projector term is the cost driver, hence Path~B; (2) refinement earns its
place where it replaces an optimizer rather than where it saves gates,
hence the filled-configuration bootstrap.}

% ===========================================================================
\section{Results}
\label{sec:results}
% ===========================================================================
\todo{To be written from the notebook.  Planned subsections and figures:}
\begin{enumerate}
  \item \textbf{Fragment spectra and path gaps.}  Low-lying seniority-zero
  spectra of $\Ni$ and $\Sn$ at each chain step; minimum gaps along Paths A
  and B versus $\mu$ and $\Lambda$.  (Fig.~2)
  \item \textbf{Refinement: overlap and energy versus $T$} for every chain
  step and both paths; $T^*$; comparison with the single-shot
  $\mathcal A_0\to\mathcal A_R$ refinement.  (Fig.~3, Table~\ref{tab:chains})
  \item \textbf{Rodeo filtering of the fragments:} $1-\mathcal F_M$ versus
  $M$ for refined inputs at $T^*$ and at a longer $T$; energy scans as a
  self-check.  (Fig.~4)
  \item \textbf{Trotterized collision:} $P_\kappa(D)$ and $S_{T|P}(D)$ versus
  step size, laboratory versus interaction picture, symmetric-subspace
  leakage, pulse truncation; overlay on Fig.~2 of
  Ref.~\cite{TwoFragment2026}.  (Fig.~5)
  \item \textbf{Exit-channel spectroscopy:} the $(E_T,E_P)$ scan at
  $D=13.4669$~fm; exclusive g.s.-to-g.s.\ probability versus $D$ against
  the inclusive $P_{+1}/P_0$; the $^{42}$Ca $0^+$ strength for $\Zr+\Ca$.
  (Fig.~6)
  \item \textbf{Entropy bracket} $\Snum\le S_{T|P}\le\Sdiag$ versus $D$ for
  both reactions.  (Fig.~7)
  \item \textbf{Gate budgets:} CNOTs for preparation (per fragment, per
  path), collision (per Trotter step and total) and readout (per accepted
  shot, per channel); total per accepted exclusive event on 30 qubits;
  verification of every stage against exact evolution.  (Fig.~8, Table~III)
  \item \textbf{Sufficient conditions versus practice} for the refinement
  paths and the adiabatic bound on $1-P_{000}$ for the pulse.
\end{enumerate}

% ===========================================================================
\section{Discussion and outlook}
\label{sec:discussion}
% ===========================================================================
\todo{Points to make: the pipeline is a single Trotterized $XY$ evolution
with three coefficient schedules and one ancilla; what the two-local
encoding buys relative to the JW pairing model; the price of 30 qubits and
of the rare-channel shot overhead; where refinement could lose (long $T^*$
if the first projectile step has a low overlap or a small gap --- here a
longer chain or a pUCCD warm start are the fallbacks, and ``optimized
refinement'' in the sense of the nucleation paper~\cite{Sadia2026} is the
variational one); extensions listed in Ref.~\cite{TwoFragment2026} ---
larger spaces, broken pairs, orbital-dependent couplings --- and which of
them the present primitives survive (broken pairs: yes, cf.\ the
particle-hole section of Ref.~\cite{EigenQPairing2026}, at the price of
leaving the two-local pair-slot encoding); hardware: the 11-qubit $\Ni$
preparation and the 26-qubit $\Zr+\Ca$ collision as device targets.}

\begin{acknowledgments}
\todo{Funding, QONTROL, FRIB, EIGEN-Q/QuNorth, computing.}
\end{acknowledgments}

\section*{Code availability}
\todo{Repository, notebook, tests.}

% ===========================================================================
\begin{thebibliography}{99}
\bibitem{TwoFragment2026} \todo{Authors}, ``Two-fragment pairing model for
sub-barrier neutron pair transfer reactions between superfluid nuclei,''
draft (2026), \texttt{paper2/twontransfer.pdf}.
\bibitem{EigenQPairing2026} S.~Bogner, C.~Glittum, H.~Hergert,
M.~Hjorth-Jensen, G.~Lange, R.~LaRose, D.~Lee, F.~P.~Massel, and J.~P.~Sadia,
``Resolution refinement followed by the rodeo algorithm for eigenstate
preparation: implementation and assessment on the pairing Hamiltonian,''
submitted to Phys.\ Rev.\ C (2026).
\bibitem{Bogner2026} \todo{resolution-refinement reference as cited in
paper 1}.
\bibitem{Choi2021} K.~Choi, D.~Lee, J.~Bonitati, Z.~Qian, and J.~Watkins,
Phys.\ Rev.\ Lett.\ \textbf{127}, 040505 (2021).
\bibitem{Qian2024} Z.~Qian \emph{et al.}, Eur.\ Phys.\ J.\ A \textbf{60},
151 (2024).
\bibitem{Jansen2007} S.~Jansen, M.-B.~Ruskai, and R.~Seiler, J.\ Math.\
Phys.\ \textbf{48}, 102111 (2007).
\bibitem{BeeLindgren2022} M.~Bee-Lindgren \emph{et al.}, Eur.\ Phys.\ J.\ A
\textbf{61}, 263 (2025).
\bibitem{Sadia2026} J.~P.~Sadia, M.~Hjorth-Jensen, D.~Lee, and R.~LaRose,
arXiv:2609.37527 (2026).
\bibitem{Corradi2011} L.~Corradi \emph{et al.}, Phys.\ Rev.\ C \textbf{84},
034603 (2011).
\bibitem{Montanari2014} D.~Montanari \emph{et al.}, Phys.\ Rev.\ Lett.\
\textbf{113}, 052501 (2014).
\bibitem{Potel2013} G.~Potel, A.~Idini, F.~Barranco, E.~Vigezzi, and
R.~A.~Broglia, Rep.\ Prog.\ Phys.\ \textbf{76}, 106301 (2013).
\end{thebibliography}

\end{document}
